% Standalone argument for review; not yet inserted into the submission.
\paragraph{Exact piecewise profile construction.}
For one side let $h$ be the cached-positive count, $n$ the unknown count,
$l=\max(\ell,h)$, and $u=\min(v,h+n)$, where $v$ is the supplied upper
count endpoint. For a discard count $d$, the two frontier witnesses are
$(a,b_0)$ and $(a,b_1)$, where
\[
 a=\min(d,u-h),\qquad b_0=\max(h,l-d),\qquad b_1=\max(h,u-d).
\]
These expressions equal the original clipped-total frontier: its first
total is $\min(u,\max(l,h+d))$ and its second total is $u$.

Fix total discards $g$ and write $x=d_L$, $d_R=g-x$. With required old
discards $o_L,o_R$, the feasible integer interval is
$[\max(o_L,g-n_R),\min(n_L,g-o_R)]$. The frontier expressions change
affine branch only at $l_L-h_L$, $u_L-h_L$,
$g-(l_R-h_R)$, and $g-(u_R-h_R)$. These four integer breakpoints
partition the interval into at most five pieces. On each closed piece
the four deficit values are integer-coefficient quadratics
\[
 Q_j(x)=(q-p)(a_L+b_L)(a_R+b_R)-q b_Lb_R,
 \qquad H_k(x,g-x)=\max_{j\leq4}Q_j(x).
\]
Branch formulas agree at shared endpoints.

\begin{theorem}[Linear-size exact profile search]
For feasible uncoupled side/key count evidence, the complete local profile,
including the lexicographically first attaining count pair at every feasible
$g$, can be constructed with $O(n_L+n_R+1)$ constant-size candidate searches.
The result includes forced old-discard lower bounds.
\end{theorem}
\begin{proof}
On one piece, include its endpoints and both neighbouring integers of every
real root of $Q_i-Q_j$ and every stationary point of each $Q_j$.
There are at most twelve pair-intersection roots and four stationary points,
so at most 34 integer candidates before duplicate removal. Between successive
real boundaries, the maximizing polynomial is fixed up to identical ties,
and its derivative has constant sign or is identically zero. The upper
envelope is therefore monotone or constant there. Its minimum over the
integers in that interval is attained at an integer nearest an endpoint;
the earliest attaining integer is included as well. Taking the smallest
pair $(H_k(x,g-x),x)$ over all candidates on all pieces gives the exact
profile value and the original left-first witness. There are
$n_L+n_R-o_L-o_R+1$ values of $g$.
\end{proof}

Root neighbours are computed exactly. For $Ax^2+Bx+C=0$, normalize $A>0$
and let $D=B^2-4AC$. When $D\geq0$, set $s=\lfloor\sqrt D\rfloor$
using integer square root and $t=s+\mathbf1\{s^2\ne D\}$. The floors of
the two roots are $\lfloor(-B-t)/(2A)\rfloor$ and
$\lfloor(-B+s)/(2A)\rfloor$. A positive integer denominator makes these
replacements preserve the quotient floors. Linear degeneracies use rational
integer division; constant differences have no isolated roots.

Let $b=O(\log(q+1)+\log(N+1))$ bound the operand bit length after capacity
clipping. The linear search count retains integer multiplication, division
and square-root costs: a bit-cost expression is
$O((n_L+n_R+1)(M(b)+D(b)+I(b)))$, with constant-factor operand expansions.
Stored profile values and witnesses use linear space. Across key-disjoint
blocks, the number of candidate searches is $O(K+N)$ per cache snapshot.
Greedy allocation then takes $O(N\log(K+1))$ heap operations; exact global
DP retains its quadratic convolution cost. Replanning costs accumulate over
snapshots. These operation bounds do not imply a speedup over vectorized
enumeration at every finite input size.
